\subsection{Healthy Serving: Small Aggregate Changes, Unequal Task Effects}
\label{sec:healthy-results}
\begin{figure*}[tp]
\centering
\includegraphics[width=\textwidth]{figures/no-failure-cost.pdf}
\caption{Failure-free protection cost across five shared admission seeds.
(a) Mean paired workflow completion difference, AegisKV minus Store-only,
for all workflows, the initial 16, and the remaining 144.
(b) Relative makespan and goodput change. Circles show seeds; squares and
whiskers show the paired seed mean and 95\% Student-$t$ interval ($df=4$).
(c) Submitted bytes by policy and destination tier; hatched portions are
warmup submissions, and crosses mark aborted bytes; nonzero aborted GiB
are labeled explicitly.
(d) Additional submitted bytes divided by Store-only's makespan versus the
later cohort's mean completion difference. Labels are seed IDs; no regression
or causal effect of traffic is inferred.}
\Description{Four panels showing paired completion differences, whole-run
changes, tier-specific protection traffic, and its association with completion.}
\label{fig:no-failure}
\end{figure*}


Figure~\ref{fig:no-failure} separates completion effects from protection
traffic. Across the five paired seeds, the all-workflow mean WCT difference
is $-0.137$\,s, with a 95\% interval of $[-0.297,0.024]$\,s.
The corresponding mean makespan change is $+0.008$\%, with an interval of
$[-0.188,0.205]$\%; individual seeds range from $-0.176$\% to $+0.201$\%.
Thus, these observations do not resolve a consistent average completion
benefit or penalty. They also do not constitute an equivalence test. The
Store-only batch completes in 329.7--331.1\,s, and its mean workflow WCT
within each seed is 32.3--32.7\,s. Goodput changes track the nearly unchanged
batch completion time because every execution emits the same prescribed
output population.

The admission cohorts show why an aggregate alone is insufficient.
The initial 16 workflows have a mean paired difference of $-0.370$\,s
($[-2.202,1.462]$\,s), whereas the later 144 have $-0.111$\,s
($[-0.413,0.191]$\,s). More pronounced redistribution occurs within each
cohort. Of the 800 paired task-runs, 340 finish sooner with AegisKV, 156
are unchanged, and 304 finish later. Paired differences
range from $-53.741$ to $+48.342$\,s. In the largest adverse case, workflow
\texttt{db12f9be559698b0} under seed 3 rises from 34.388 to 82.731\,s.
Small mean overhead therefore coexists with substantial task-level harm.

Protection has a substantial byte cost despite the small mean makespan
change. Store-only submits 2,825.4\,GiB of protection traffic per run.
AegisKV submits 1,496.5--1,586.7\,GiB to Peer-HBM and
3,133.0--3,183.8\,GiB to Store-DRAM. Across seeds, the additional submitted
traffic averages approximately 1,874\,GiB, or 66.3\% of the comparator's
traffic. These are cumulative copy bytes, not the instantaneous protected
KV footprint. The two tiers may protect overlapping prompt intervals.
Aborts account for 0.5625 and 3.9375\,GiB of Peer submissions in seeds 2
and 4, respectively, and zero bytes in the other cases. Submitted traffic
and bytes actually transferred are therefore kept separate.

The incremental submitted-byte rate is 5.62--5.75\,GiB/s when divided by
Store-only's batch makespan. Figure~\ref{fig:no-failure}(d) associates those
five rates with later-cohort WCT changes of both signs. It does not identify
a causal response to bandwidth: admission order, cache reuse, placement,
and resource contention change together. These costs and task-level effects are specific to the declared operating point.

\paragraph{Tracing the largest adverse task.}
The complete scheduler record explains the 48.342\,s difference without
attributing it to a cache miss. For this 27-call task, the final fan-in call
is released 16.961\,s later relative to workflow admission, and its own
release-to-completion interval increases by 31.381\,s. Both policies reuse
the same 64-token normal-HBM prefix. The call actually waits less before
its first scheduler step with AegisKV: 2.258 versus 4.761\,s. However,
the time in steps containing that call grows from 8.082 to 41.966\,s,
with no gaps between its participating steps in either execution.
Its summed ideal participant shares are identical at 0.764\,s. The
shared batches differ: 108 mixed prefill/decode steps with AegisKV versus
20 with Store-only. Reconstructing each step's maximum participant share
recovers the recorded durations exactly. This identifies shared-step
composition as the immediate source of the longer executing interval
under this timing model; it does not isolate a causal effect of protection
traffic from the accompanying changes in routing and batching.
